\documentclass[10pt,a4paper]{amsart}
\usepackage{amsmath,amssymb,amsthm,graphicx}
\usepackage{verbatim}
\newenvironment{reference}{\comment}{\endcomment}
%\newenvironment{reference}{}{}
\newenvironment{slogan}{\comment}{\endcomment}
\newenvironment{history}{\comment}{\endcomment}

% For commutative diagrams we use Xy-pic
\usepackage[all]{xy}

% We use 2cell for 2-commutative diagrams.
\xyoption{2cell}
% Theorem environments.
%
\theoremstyle{plain}
\newtheorem{theorem}[subsection]{Theorem}
\newtheorem{proposition}[subsection]{Proposition}
\newtheorem{lemma}[subsection]{Lemma}

\theoremstyle{definition}
\newtheorem{definition}[subsection]{Definition}
\newtheorem{example}[subsection]{Example}
\newtheorem{exercise}[subsection]{Exercise}
\newtheorem{situation}[subsection]{Situation}

\theoremstyle{remark}
\newtheorem{remark}[subsection]{Remark}
\newtheorem{remarks}[subsection]{Remarks}

\numberwithin{equation}{subsection}

% Macros
%
\def\lim{\mathop{\mathrm{lim}}\nolimits}
\def\colim{\mathop{\mathrm{colim}}\nolimits}
\def\Spec{\mathop{\mathrm{Spec}}}
\def\Hom{\mathop{\mathrm{Hom}}\nolimits}
\def\Ext{\mathop{\mathrm{Ext}}\nolimits}
\def\SheafHom{\mathop{\mathcal{H}\!\mathit{om}}\nolimits}
\def\SheafExt{\mathop{\mathcal{E}\!\mathit{xt}}\nolimits}
\def\Sch{\mathit{Sch}}
\def\Mor{\mathop{\mathrm{Mor}}\nolimits}
\def\Ob{\mathop{\mathrm{Ob}}\nolimits}
\def\Sh{\mathop{\mathit{Sh}}\nolimits}
\def\NL{\mathop{N\!L}\nolimits}
\def\CH{\mathop{\mathrm{CH}}\nolimits}
\def\proetale{{pro\text{-}\acute{e}tale}}
\def\etale{{\acute{e}tale}}
\def\QCoh{\mathit{QCoh}}
\def\Ker{\mathop{\mathrm{Ker}}}
\def\Im{\mathop{\mathrm{Im}}}
\def\Coker{\mathop{\mathrm{Coker}}}
\def\Coim{\mathop{\mathrm{Coim}}}

% Boxtimes
%
\DeclareMathSymbol{\boxtimes}{\mathbin}{AMSa}{"02}

%
% Macros for moduli stacks/spaces
%
\def\QCohstack{\mathcal{QC}\!\mathit{oh}}
\def\Cohstack{\mathcal{C}\!\mathit{oh}}
\def\Spacesstack{\mathcal{S}\!\mathit{paces}}
\def\Quotfunctor{\mathrm{Quot}}
\def\Hilbfunctor{\mathrm{Hilb}}
\def\Curvesstack{\mathcal{C}\!\mathit{urves}}
\def\Polarizedstack{\mathcal{P}\!\mathit{olarized}}
\def\Complexesstack{\mathcal{C}\!\mathit{omplexes}}
% \Pic is the operator that assigns to X its picard group, usage \Pic(X)
% \Picardstack_{X/B} denotes the Picard stack of X over B
% \Picardfunctor_{X/B} denotes the Picard functor of X over B
\def\Pic{\mathop{\mathrm{Pic}}\nolimits}
\def\Picardstack{\mathcal{P}\!\mathit{ic}}
\def\Picardfunctor{\mathrm{Pic}}
\def\Deformationcategory{\mathcal{D}\!\mathit{ef}}
\usepackage{cleveref}
\crefname{originalequation}{original equation}{original equations}
\Crefname{originalequation}{Original equation}{Original equations}
\makeatletter
\newlabel{equation-base-change-map}{{0AA6}{source}{}{sourcepointer-1}{}}
\newlabel{equation-base-change-map@cref}{{[originalequation][0][]0AA6}{[1][0][]source}{}{}{}}
\makeatother
\crefname{originallemma}{original lemma}{original lemmas}
\Crefname{originallemma}{Original lemma}{Original lemmas}
\makeatletter
\newlabel{lemma-twisted-inverse-image}{{0A9E}{source}{}{sourcepointer-2}{}}
\newlabel{lemma-twisted-inverse-image@cref}{{[originallemma][0][]0A9E}{[1][0][]source}{}{}{}}
\makeatother
\crefname{originallemma}{original lemma}{original lemmas}
\Crefname{originallemma}{Original lemma}{Original lemmas}
\makeatletter
\newlabel{categories-lemma-properties-2-cat-cats}{{003F}{source}{}{sourcepointer-3}{}}
\newlabel{categories-lemma-properties-2-cat-cats@cref}{{[originallemma][0][]003F}{[1][0][]source}{}{}{}}
\makeatother
\crefname{originaldefinition}{original definition}{original definitions}
\Crefname{originaldefinition}{Original definition}{Original definitions}
\makeatletter
\newlabel{categories-definition-horizontal-composition}{{003E}{source}{}{sourcepointer-4}{}}
\newlabel{categories-definition-horizontal-composition@cref}{{[originaldefinition][0][]003E}{[1][0][]source}{}{}{}}
\makeatother
\crefname{originalreference}{original reference}{original references}
\Crefname{originalreference}{Original reference}{Original references}
\makeatletter
\newlabel{cohomology-remark-compose-base-change}{{0ATL}{source}{}{sourcepointer-5}{}}
\newlabel{cohomology-remark-compose-base-change@cref}{{[originalreference][0][]0ATL}{[1][0][]source}{}{}{}}
\makeatother
\title{The base change maps of two stacked cartesian squares compose to the base change map of the outer rectangle}
\author{The Stacks Project Authors}
\date{stacks-project:duality.tex:local-snapshot-20261012; GNU FDL 1.2 or later}
\begin{document}
\maketitle
\noindent\textit{Editorial note.} All delivered mathematical text is the author's own, copied from the pinned Stacks Project chapter file duality.tex; no statement or proof step was rewritten, reordered or removed.\par\medskip
\noindent\textit{Editorial note.} Preamble provenance: the pinned chapter file opens with an input of the shared preamble and carries no preamble text of its own, so the project's single Stacks-wide preamble staged at sources/stacks-duality/preamble.tex is excerpted verbatim under the recipe fixed by LANE\_BRIEF 22.1, as three spans, lines 8-13 (verbatim package and the reference/slogan/history environments), lines 14-18 (Xy-pic setup) and lines 154-217 (the native theorem declarations and mathematical macros). Nothing was added, removed or redefined; the wrapper class and the amsmath/amssymb/amsthm/graphicx base packages are supplied by the assembler.\par\medskip
\noindent\textit{Editorial note.} The counted statement and proof contain no cross reference, no citation and no included file, so the delivered body is self contained.\par\medskip
\noindent\textit{Editorial note.} Canonical location of the counted proposition: Stacks tag 0ATQ (sources/stacks-duality/pinned-original-source.tex line 1238).\par\medskip
\noindent\textit{Allowed standard prerequisites (editorial).}\begin{itemize}
\item The base change map between the right adjoints for a cartesian Tor independent square. Same chapter, equation labelled equation-base-change-map and Lemma labelled lemma-twisted-inverse-image, cited at proof lines 1268-1276. Supplies the three arrows whose composition is asserted.
\item The dual statement for ordinary base change maps. Cohomology, Remark labelled cohomology-remark-compose-base-change, cited at proof lines 1368-1370. Identifies the final reduced arrow with xi\_\{rect\}.
\item Coherence of the compositions of functors in a two category. Categories, Lemma labelled categories-lemma-properties-2-cat-cats and Definition labelled categories-definition-horizontal-composition, cited at proof lines 1322-1326 and 1350-1352. Gives the commutativity of the quadrilaterals and the triangle used in the chase.
\end{itemize}\medskip
\ifdefined\hypertarget\hypertarget{sourcepointer-1}{}\fi
\noindent\textit{Source reference.} Original equation 0AA6: sources/stacks-duality/pinned-original-source.tex, source lines 1192--1192. Same-chapter pointer into the author environment equation at source line 1192; invoked at proof line None of the counted proof. It is retained only as a navigation pointer; every load-bearing step is inside the delivered span.\par\medskip
\ifdefined\hypertarget\hypertarget{sourcepointer-2}{}\fi
\noindent\textit{Source reference.} Original lemma 0A9E: sources/stacks-duality/pinned-original-source.tex, source lines 653--653. Same-chapter pointer into the author environment lemma at source line 653; invoked at proof line 1260 of the counted proof. It is retained only as a navigation pointer; every load-bearing step is inside the delivered span.\par\medskip
\ifdefined\hypertarget\hypertarget{sourcepointer-3}{}\fi
\noindent\textit{Source reference.} Original lemma 003F: sources/stacks-chow/tags.txt, source lines 132--132. Cross-chapter standard prerequisite located in the project tag list staged in the local snapshot; invoked at proof line 1316 of the counted proof. It is retained only as a navigation pointer; every load-bearing step is inside the delivered span.\par\medskip
\ifdefined\hypertarget\hypertarget{sourcepointer-4}{}\fi
\noindent\textit{Source reference.} Original definition 003E: sources/stacks-chow/tags.txt, source lines 131--131. Cross-chapter standard prerequisite located in the project tag list staged in the local snapshot; invoked at proof line 1318 of the counted proof. It is retained only as a navigation pointer; every load-bearing step is inside the delivered span.\par\medskip
\ifdefined\hypertarget\hypertarget{sourcepointer-5}{}\fi
\noindent\textit{Source reference.} Original reference 0ATL: sources/stacks-chow/tags.txt, source lines 13619--13619. Cross-chapter standard prerequisite located in the project tag list staged in the local snapshot; invoked at proof line 1370 of the counted proof. It is retained only as a navigation pointer; every load-bearing step is inside the delivered span.\par\medskip
\begin{lemma}
\label{lemma-compose-base-change-maps}
Consider a commutative diagram
$$
\xymatrix{
X' \ar[r]_k \ar[d]_{f'} & X \ar[d]^f \\
Y' \ar[r]^l \ar[d]_{g'} & Y \ar[d]^g \\
Z' \ar[r]^m & Z
}
$$
of quasi-compact and quasi-separated schemes where
both diagrams are cartesian and where $f$ and $l$
as well as $g$ and $m$ are Tor independent.
Then the maps (\ref{equation-base-change-map})
for the two squares compose to give the base
change map for the outer rectangle (see proof for a precise statement).
\end{lemma}

\begin{proof}
It follows from the assumptions that $g \circ f$ and $m$ are Tor
independent (details omitted), hence the statement makes sense.
In this proof we write $k^*$ in place of $Lk^*$ and $f_*$ instead
of $Rf_*$. Let $a$, $b$, and $c$ be the right adjoints of
Lemma \ref{lemma-twisted-inverse-image}
for $f$, $g$, and $g \circ f$ and similarly for the primed versions.
The arrow corresponding to the top square is the composition
$$
\gamma_{top} :
k^* \circ a \to k^* \circ a \circ l_* \circ l^*
\xleftarrow{\xi_{top}} k^* \circ k_* \circ a' \circ l^* \to a' \circ l^*
$$
where $\xi_{top} : k_* \circ a' \to a \circ l_*$
is an isomorphism (hence can be inverted)
and is the arrow ``dual'' to the base change map
$l^* \circ f_* \to f'_* \circ k^*$. The outer arrows come
from the canonical maps $1 \to l_* \circ l^*$ and $k^* \circ k_* \to 1$.
Similarly for the second square we have
$$
\gamma_{bot} :
l^* \circ b \to l^* \circ b \circ m_* \circ m^*
\xleftarrow{\xi_{bot}} l^* \circ l_* \circ b' \circ m^* \to b' \circ m^*
$$
For the outer rectangle we get
$$
\gamma_{rect} :
k^* \circ c \to k^* \circ c \circ m_* \circ m^*
\xleftarrow{\xi_{rect}} k^* \circ k_* \circ c' \circ m^* \to c' \circ m^*
$$
We have $(g \circ f)_* = g_* \circ f_*$ and hence
$c = a \circ b$ and similarly $c' = a' \circ b'$.
The statement of the lemma is that $\gamma_{rect}$
is equal to the composition
$$
k^* \circ c = k^* \circ a \circ b \xrightarrow{\gamma_{top}}
a' \circ l^* \circ b \xrightarrow{\gamma_{bot}}
a' \circ b' \circ m^* = c' \circ m^*
$$
To see this we contemplate the following diagram:
$$
\xymatrix{
& & k^* \circ a \circ b \ar[d] \ar[lldd] \\
& & k^* \circ a \circ l_* \circ l^* \circ b \ar[ld] \\
k^* \circ a \circ b \circ m_* \circ m^* \ar[r] &
k^* \circ a \circ l_* \circ l^* \circ b \circ m_* \circ m^* &
k^* \circ k_* \circ a' \circ l^* \circ b \ar[u]_{\xi_{top}} \ar[d] \ar[ld] \\
& k^*\circ k_* \circ a' \circ l^* \circ b \circ m_* \circ m^*
\ar[u]_{\xi_{top}} \ar[rd] &
a' \circ l^* \circ b \ar[d] \\
k^* \circ k_* \circ a' \circ b' \circ m^* \ar[uu]_{\xi_{rect}} \ar[ddrr] &
k^*\circ k_* \circ a' \circ l^* \circ l_* \circ b' \circ m^*
\ar[u]_{\xi_{bot}} \ar[l] \ar[dr] &
a' \circ l^* \circ b \circ m_* \circ m^* \\
& & a' \circ l^* \circ l_* \circ b' \circ m^* \ar[u]_{\xi_{bot}} \ar[d] \\
& & a' \circ b' \circ m^*
}
$$
Going down the right hand side we have the composition and going
down the left hand side we have $\gamma_{rect}$.
All the quadrilaterals on the right hand side of this diagram commute
by Categories, Lemma \ref{categories-lemma-properties-2-cat-cats}
or more simply the discussion preceding
Categories, Definition \ref{categories-definition-horizontal-composition}.
Hence we see that it suffices to show the diagram
$$
\xymatrix{
a \circ l_* \circ l^* \circ b \circ m_* &
a \circ b \circ m_* \ar[l] \\
k_* \circ a' \circ l^* \circ b \circ m_* \ar[u]_{\xi_{top}} & \\
k_* \circ a' \circ l^* \circ l_* \circ b' \ar[u]_{\xi_{bot}} \ar[r] &
k_* \circ a' \circ b' \ar[uu]_{\xi_{rect}}
}
$$
becomes commutative if we invert the arrows $\xi_{top}$, $\xi_{bot}$,
and $\xi_{rect}$ (note that this is different from asking the
diagram to be commutative). However, the diagram
$$
\xymatrix{
& a \circ l_* \circ l^* \circ b \circ m_* \\
a \circ l_* \circ l^* \circ l_* \circ b'
\ar[ru]^{\xi_{bot}} & &
k_* \circ a' \circ l^* \circ b \circ m_* \ar[ul]_{\xi_{top}} \\
& k_* \circ a' \circ l^* \circ l_* \circ b'
\ar[ul]^{\xi_{top}} \ar[ur]_{\xi_{bot}}
}
$$
commutes by Categories, Lemma \ref{categories-lemma-properties-2-cat-cats}.
Since the diagrams
$$
\vcenter{
\xymatrix{
a \circ l_* \circ l^* \circ b \circ m_* & a \circ b \circ m \ar[l] \\
a \circ l_* \circ l^* \circ l_* \circ b' \ar[u] &
a \circ l_* \circ b' \ar[l] \ar[u]
}
}
\quad\text{and}\quad
\vcenter{
\xymatrix{
a \circ l_* \circ l^* \circ l_* \circ b' \ar[r] & a \circ l_* \circ b' \\
k_* \circ a' \circ l^* \circ l_* \circ b' \ar[u] \ar[r] &
k_* \circ a' \circ b' \ar[u]
}
}
$$
commute (see references cited) and since the composition of
$l_* \to l_* \circ l^* \circ l_* \to l_*$ is the identity,
we find that it suffices to prove that
$$
k \circ a' \circ b' \xrightarrow{\xi_{bot}} a \circ l_* \circ b
\xrightarrow{\xi_{top}} a \circ b \circ m_*
$$
is equal to $\xi_{rect}$ (via the identifications $a \circ b = c$
and $a' \circ b' = c'$). This is the statement dual to
Cohomology, Remark \ref{cohomology-remark-compose-base-change}
and the proof is complete.
\end{proof}

\end{document}
